\section{Part 1 --- Hormonal Control of Glucose}

This part implements the Cobelli \emph{et~al.}\ (1982) intermediate-complexity model of
glucose--insulin--glucagon regulation, as presented in Haefner Ch.~12, and works through
exercises 1--7 of \S12.5. The core is written in Python; integration uses
\texttt{scipy.integrate.solve\_ivp} (LSODA), with a Radau cross-check. All exercises run against a
single verified core (\texttt{cobelli.py}). Before the exercises I describe the model and the one
non-trivial step needed to make it run at all: a small, documented calibration of an under-specified
parameter table.

\subsection{The Cobelli intermediate-complexity model}

The model tracks seven state variables in three coupled subsystems, structured as a Forrester
negative-feedback diagram (Fig.~12.2 of the text):

\begin{itemize}[leftmargin=1.4em,itemsep=1pt]
\item \textbf{Glucose subsystem:} $g$ = glucose in plasma and extracellular fluid (mg).
\item \textbf{Glucagon subsystem:} $c$ = glucagon in plasma and interstitial fluid (nU).
\item \textbf{Insulin subsystem (five compartments):} $s$ = stored pancreatic insulin, $r$ =
promptly releasable pancreatic insulin, $l$ = liver insulin, $p$ = plasma insulin, $i$ =
interstitial-fluid insulin (all \si{\micro U}).
\end{itemize}

Each state variable has an associated concentration $\bar{g}=g/V_b$, $\bar{p}=p/V_p$, and so on,
where the volumes $V_b,V_p,V_l,V_i$ are fixed fractions of body weight. Feedback is expressed through
fluxes built from the hyperbolic tangent: a stimulatory effect uses
$0.5\,[1+\tanh(b(\Delta\bar{x}+c))]$ and an inhibitory effect uses $0.5\,[1-\tanh(b(\Delta\bar{x}+c))]$,
where $\Delta\bar{x}$ is the deviation of a concentration from its basal value, $b$ sets the slope and
$c$ the centre. Because $\tanh$ has range $\pm1$, each gate is bounded in $[0,1]$ and saturates
smoothly; multiplying several gates implements ``multiple limiting factors''. The governing equations
(Eqs.~12.1--12.7) are

\begin{align}
\dd{g}{t} &= \mathrm{NHGB} - F_3 - F_4 - F_5 + I_g(t), \\
\dd{c}{t} &= -h_{02}\,c + F_7, \\
\dd{i}{t} &= -m_{13}\,i + m_{31}\,p, \\
\dd{l}{t} &= -(m_{02}+m_{12})\,l + m_{21}\,p + F_6, \\
\dd{p}{t} &= -(m_{01}+m_{21}+m_{31})\,p + m_{12}\,l + m_{13}\,i + I_p(t), \\
\dd{r}{t} &= k_{21}\,s - k_{12}\,r - F_6, \\
\dd{s}{t} &= -k_{21}\,s + k_{12}\,r + W,
\end{align}

where $\mathrm{NHGB}=F_1-F_2$ is the net hepatic glucose balance (liver production minus uptake),
$F_3$ is renal excretion, $F_4$ peripheral (muscle/adipose) use, $F_5$ the nearly
insulin-independent CNS/red-blood-cell sink, $F_6$ pancreatic insulin secretion, $F_7$ glucagon
secretion, $W$ insulin synthesis into the stored pool, and $I_g,I_p$ are the glucose/insulin input
functions used to drive diagnostic tests. The secretion flux
$F_6=0.5\,a_6[1+\tanh(b_6(\Delta\bar{g}+c_6))]\,r$ is proportional to the releasable pool $r$, which
matters for the ``hump'' in Exercise~2.

\textbf{Transcription fix.} The renal-excretion gate is printed in the book as
$M_{31}=0.5[1+\tanh(b_{13}(\bar{g}+c_{13}))]$, but $b_{13}=2.75\times10^{-2}$, $c_{13}=20$ are glucose
self-suppression constants for $M_1$. Table~12.2 instead lists $b_{31}=20$, $c_{31}=-180$, which set a
renal threshold near $180\mgdl$ --- the physiologically correct reading. We therefore use
$b_{31},c_{31}$ (so $F_3$ stays near zero until glucose exceeds the threshold), treating the printed
$b_{13}$ as a typo.

\subsection{Calibration: reproducing Fig.~12.4 from an under-specified table}
\label{sec:calibration}

Taken literally and in the amount-based form of Eqs.~12.1--12.7, the Table~12.2 parameters do
\emph{not} reproduce Figs.~12.4a/b. This was verified independently and is a genuine gap in the
printed table, not a units error on our side (an amount$\leftrightarrow$concentration $V_b\approx140$
factor closes less than 1\% of it). Quantitatively, for a 70\,kg subject:

\begin{itemize}[leftmargin=1.4em,itemsep=1pt]
\item The literal basal glucose clearance is $F_3+F_4+F_5\approx8.4\times10^{-4}\,$mg/min against a
basal glucose amount of $12{,}810$\,mg, giving a clearance timescale of $\sim1.5\times10^{7}$\,min.
Glucose essentially never clears (Fig.~12.4 shows recovery in $\sim$90\,min).
\item The basal plasma insulin amount is $34{,}650\,\si{\micro U}$, which needs a basal secretion of
$F_6\approx1.6\times10^{4}\,\si{\micro U}$/min, but the literal synthesis flux is
$W(\text{basal})\approx0.017\,\si{\micro U}$/min --- about $10^{6}$ times too small. The insulin and
glucagon axes collapse to zero and no IVGTT response is possible.
\end{itemize}

The model is explicitly described as ``semi-phenomenological\,\dots\ tailored/scaled to a particular
patient'' (Haefner \S12.3). Sitting the model at a patient's basal operating point by choosing the
under-determined source/clearance scales is exactly its intended use, so we applied a minimal,
structure-preserving calibration and changed \emph{no} feedback shape:

\begin{enumerate}[leftmargin=1.5em,itemsep=2pt]
\item A uniform \texttt{glucose\_flux\_scale} (default $10^{5}$) multiplies the four glucose fluxes
(NHGB, $F_3$, $F_4$, $F_5$), setting the clearance timescale to $\sim$90\,min. Because
$\mathrm{NHGB}=F_3+F_4+F_5$ at basal, a \emph{common} scale cancels, so the basal steady state and
every $\tanh$ feedback shape are exactly unchanged (the basal residual stays at $3.6\times10^{-12}$
for scale $\in\{1,10^{3},10^{5},10^{7}\}$; setting it to $1.0$ recovers the literal table).
\item The source scales $a_w$ (insulin synthesis) and $a_{71}$ (glucagon secretion) are chosen so the
published clinical basal is an exact steady state. No $\tanh$ slope $b$, centre $c$, or transfer rate
($m_{ij}$, $k_{ij}$) is altered.
\item $b_{52}$, a constant offset in the $F_5$ CNS/RBC sink, is shifted by the tiny published
NHGB$-(F_3+F_4+F_5)$ imbalance so that $\dd{g}{t}=0$ at basal exactly. This is a constant offset,
not a shape change.
\end{enumerate}

After calibration the basal residual is $3.6\times10^{-12}$; a normal IVGTT gives a glucose peak of
$250\mgdl$ recovering in $\sim$52\,min, plasma insulin rising $11\to40\uU$ then decaying, and glucagon
dipping then recovering, with LSODA and Radau agreeing to $1.6\times10^{-4}$. These match Fig.~12.4
qualitatively (Fig.~\ref{fig:verify}).

\textbf{What is not claimed.} We do not claim that $10^{5}$, or the recalibrated $a_w,a_{71},b_{52}$,
are the original Cobelli (1982) coefficients. They are a documented fit that reproduces the published
figures from an under-specified table while leaving all feedback structure intact.

\begin{figure}[t]
\centering
\includegraphics[width=0.62\linewidth]{verify_normal_ivgtt.png}
\caption{Verification of the calibrated core against a normal IVGTT. Plasma glucose peaks near
$250\mgdl$ and recovers within $\sim$52\,min; plasma insulin rises and decays; LSODA and Radau agree
to $1.6\times10^{-4}$. The calibration preserves the basal steady state and all $\tanh$ feedback
shapes.}
\label{fig:verify}
\end{figure}

\subsection{Exercise 1: Normal IVGTT and where glucose goes}

\textbf{Question.} Code the model, reproduce Figs.~12.4a/b (plasma glucose; plasma/liver/interstitial
insulin), additionally plot glucagon and the liver uptake $F_2$, and answer: where does the majority
of glucose go?

\textbf{Approach.} A normal 70\,kg subject receives a single IV glucose bolus of $0.33\,$g/kg
($=23.1$\,g) over 2\,min, integrated from the calibrated basal state for 150\,min. The glucose-disposal
fluxes are integrated over the window to apportion the load.

Plasma glucose rises from basal $91.5$ to a peak of $250\mgdl$ at $t=2$\,min, recovering within 5\% of
basal by $t=52$\,min. Plasma insulin rises from $11$ to a peak of $40\uU$ within minutes, and glucagon
is suppressed from $75$ down to $6$\,nU/ml by the combined glucose and insulin rise
(Figs.~\ref{fig:ex1_glucose}--\ref{fig:ex1_glucagon}). The cumulative disposal split is given in
Table~\ref{tab:ex1_split}.

\begin{figure}[H]
\centering
\begin{subfigure}{0.48\linewidth}
\includegraphics[width=\linewidth]{ex1_glucose.png}
\caption{Plasma glucose (Fig.~12.4a).}
\end{subfigure}\hfill
\begin{subfigure}{0.48\linewidth}
\includegraphics[width=\linewidth]{ex1_insulin.png}
\caption{Plasma / liver / interstitial insulin (Fig.~12.4b).}
\end{subfigure}
\caption{Normal IVGTT ($23.1$\,g over 2\,min). Glucose peaks at $250\mgdl$ then recovers; all three
insulin compartments rise within minutes of the bolus.}
\label{fig:ex1_glucose}
\end{figure}

\begin{table}[H]
\centering
\caption{Cumulative glucose disposal over the 150-min IVGTT response. Peripheral disposal
$F_3+F_4+F_5=40{,}170$\,mg.}
\label{tab:ex1_split}
\begin{tabular}{lrr}
\toprule
Sink & Cumulative (mg) & \% of peripheral disposal \\
\midrule
$F_5$ CNS / RBC & $24{,}503$ & $61.0$ \\
$F_4$ muscle / adipose & $11{,}202$ & $27.9$ \\
$F_3$ renal & $4{,}465$ & $11.1$ \\
\midrule
$\int\mathrm{NHGB}$ (signed) & $16{,}855$ & --- \\
\bottomrule
\end{tabular}
\end{table}

\textbf{Discussion.} The majority of disposed glucose goes to the CNS and red blood cells ($F_5$,
$61.0\%$), followed by insulin-sensitive muscle and adipose tissue ($F_4$, $27.9\%$); renal excretion
($F_3$) contributes $11.1\%$. This is physiologically sensible: $F_5$ is a large, nearly
insulin-independent sink, while $F_3$ stays near zero until $\bar{g}$ exceeds the $\sim$180\,mg/100\,ml
renal threshold --- here the peak is $250\mgdl$, so $F_3$ is engaged only modestly. In Forrester terms,
the bolus drives glucose up, which opens $F_6$ ($\propto r$) so insulin climbs, which both clears
glucose peripherally and (via inhibitory $-\tanh$ gates) suppresses glucagon. The sign of
$\int\mathrm{NHGB}$ ($+16{,}855$\,mg) indicates the liver is a net source over the full window, as the
basal hepatic output dominates once glucose has cleared.

\begin{figure}[H]
\centering
\begin{subfigure}{0.48\linewidth}
\includegraphics[width=\linewidth]{ex1_glucagon_F2.png}
\caption{Glucagon and liver uptake $F_2$.}
\end{subfigure}\hfill
\begin{subfigure}{0.48\linewidth}
\includegraphics[width=\linewidth]{ex1_disposal_split.png}
\caption{Cumulative disposal fluxes.}
\end{subfigure}
\caption{Glucagon is suppressed during the excursion; the cumulative-flux panel shows CNS/RBC ($F_5$)
as the dominant sink.}
\label{fig:ex1_glucagon}
\end{figure}

\subsection{Exercise 2: The insulin ``hump''}

\textbf{Question.} Why does the ``hump'' in insulin concentration develop in Fig.~12.5 (repeated
glucose pulses, normal subject), and why is there a similar hump for obese persons in Fig.~12.7?

\textbf{Approach.} (i) A normal subject is given an escalating repeated-pulse train
($[0.5,1.0,2.5,5.0,10.0,20.0,40.0]$\,g at $t=[0,30,70,110,180,270,400]$\,min) over 700\,min. (ii) An
obese subject is given a single $23.1$\,g IVGTT from its pre-equilibrated basal.

In the normal repeated-pulse run a hump is detected: insulin re-rises to $\sim$37\uU\ near $t=450$\,min
after dipping to $\sim$31\uU\ (Fig.~\ref{fig:ex2}). In the obese single-IVGTT run a hump is likewise
detected, with a secondary insulin rise to $\sim$32\uU\ near $t=105$\,min (dip $\sim$22\uU; peak
$36\uU$ vs basal $11\uU$).

\textbf{Mechanism.} Both humps share one origin: secretion
$F_6=0.5\,a_6[1+\tanh(b_6(\Delta\bar{g}+c_6))]\,r$ is proportional to the releasable pool $r$, which
refills from the stored pool $s$ on its own slow timescale ($\dd{r}{t}=k_{21}s-k_{12}r-F_6$;
$\dd{s}{t}=-k_{21}s+k_{12}r+W$). You cannot secrete insulin you do not have, and $r$ cannot refill
instantly. In Fig.~12.5 each pulse drains $r$; because pulses are spaced shorter than the
refill/clearance time, $r$ is repeatedly depleted and over-refilled from $s$, and the multi-compartment
lag ($r\to l\to p\to i$) stacks the delayed contributions, so plasma insulin does not decay
exponentially between pulses but accumulates into a delayed hump after the larger pulses. In the obese
case (Fig.~12.7) the same pool dynamics are driven by parameters rather than repeated input: $b_6=0.5$
(vs nominal $9.23\times10^{-2}$) steepens the $F_6$ gate so a single load triggers an outsized
drain-and-refill of $r$, and $m_{02}=0.13$ slows downstream clearance, so the pools overshoot and
secretion re-engages as $r$ refills --- an abnormal hump on the insulin-decay curve even though the
glucose response is near-normal.

\begin{figure}[H]
\centering
\begin{subfigure}{0.48\linewidth}
\includegraphics[width=\linewidth]{ex2_normal_hump.png}
\caption{Normal, repeated pulses.}
\end{subfigure}\hfill
\begin{subfigure}{0.48\linewidth}
\includegraphics[width=\linewidth]{ex2_obese_hump.png}
\caption{Obese, single IVGTT.}
\end{subfigure}
\caption{The insulin hump in both cases. Panels show glucose, insulin, the $r/s$ pools and $F_6$: the
depletion--refill cycling of $r$ drives the secondary insulin rise.}
\label{fig:ex2}
\end{figure}

\subsection{Exercise 3: Repeated pulses, diabetic vs normal}

\textbf{Question.} Apply the Fig.~12.5 pulse sequence to a diabetic subject and compare with a normal
subject.

\textbf{Approach.} The identical escalating train is applied to both over 800\,min using the single
canonical core. The diabetic uses the Table~12.3 secretion/uptake-shape changes, with the
insulin-driven hepatic ($F_2$) and peripheral ($H_4$) uptake gates frozen, and a calibration to its
own elevated operating point.

\begin{table}[H]
\centering
\caption{Repeated-pulse response, normal vs diabetic.}
\label{tab:ex3}
\begin{tabular}{lrr}
\toprule
Quantity & Normal & Diabetic \\
\midrule
Basal glucose ($\mgdl$) & $91.5$ & $140.0$ \\
Peak glucose ($\mgdl$) & $366$ & $418$ \\
Final glucose ($\mgdl$) & $89.9$ & $139.7$ \\
Basal insulin ($\uU$) & $11.0$ & $7.0$ \\
Peak insulin ($\uU$) & $62$ & $18.3$ \\
Recovery after last pulse (min) & $462$ & $524$ \\
\bottomrule
\end{tabular}
\end{table}

\textbf{Disclosed modeling choice.} Table~12.3 lists $F_2=0.037$, but in the textbook's native
(un-scaled) units. Fed through our $\texttt{glucose\_flux\_scale}=10^{5}$, the literal value becomes a
$\sim$3700\,mg/min hepatic ``vacuum'' ($\sim$70$\times$ the normal basal $F_2\approx52$\,mg/min) that
clamps glucose flat --- physiologically backwards. The scale-consistent reading of ``$F_2/H_2$ replaced
by a constant'' is that the diabetic liver can no longer \emph{ramp} uptake with insulin, so we freeze
$F_2$ near the normal basal hepatic-uptake flux and freeze $H_4$ at the Table~12.3 value $0.0012$. The
elevated basal does not emerge from freezing alone (the $M_1$ glucose self-suppression and $G_1$
glucagon gates make $91.5\mgdl$ a strong attractor), so the diabetic is calibrated to its own operating
point ($\mathrm{DIABETIC\_BASAL}=140/7/75$) using the same under-determined scales ($a_w,a_{71},b_{52}$)
as the normal subject --- no feedback shape is changed. At the low liver-insulin diabetic state the
\emph{unfrozen} hepatic-production gate $H_1$ naturally sits high ($\sim$0.66 vs $0.25$ normal):
hepatic production is no longer insulin-suppressed, the textbook diabetic mechanism, emergent rather
than forced.

\textbf{Discussion.} The diabetic basal glucose ($140\mgdl$) is elevated and its basal insulin
($7.0\uU$) blunted, both relative to normal ($91.5\mgdl$, $11.0\uU$). Its peak insulin ($18.3\uU$) is
far below the normal $62\uU$ --- the blunted Table~12.3 secretion machinery cannot mount the same
overshoot, so the characteristic insulin hump is suppressed. With impaired uptake and unsuppressed
production, the diabetic clears each pulse more slowly and stacks glucose above its already-high basal.
On a standard IVGTT this same configuration recovers in $\sim$100\,min vs $\sim$52\,min for normal
($\sim$1.9$\times$), the model expression of the textbook's ``$\sim$2$\times$ slower recovery''
(\S12.3.6).

\begin{figure}[t]
\centering
\includegraphics[width=0.66\linewidth]{ex3_diabetic_vs_normal.png}
\caption{Repeated-pulse plasma glucose and insulin, normal vs diabetic (overlay and side-by-side). The
diabetic stays elevated above its higher basal with a blunted, hump-free insulin response.}
\label{fig:ex3}
\end{figure}

\subsection{Exercise 4: Meal (60-min infusion) vs IVGTT pulse}

\textbf{Question.} Deliver glucose not as an instantaneous IVGTT pulse but as a constant input spread
over 60\,min (a meal): 25\,g to a 70\,kg subject. Compare with an IVGTT of the same total dose.

\textbf{Approach.} A normal subject receives $25$\,g either as a constant infusion
$I_g=416.7$\,mg/min over $[0,60]$\,min, or as a $\sim$2-min bolus ($I_g=12{,}500$\,mg/min). The only
difference between runs is the temporal spread of the input. The 60-min rectangular infusion is a
deliberate first-order stand-in for the smooth, delayed gastrointestinal rate of appearance $R_a(t)$ of
a real meal (cf.\ Dalla Man \emph{et~al.}\ 2007); the full GI submodel is not required here.

\begin{table}[H]
\centering
\caption{Meal (60-min infusion) versus IVGTT (2-min bolus), same 25\,g total. Basal glucose
$91.5\mgdl$, insulin $11.0\uU$.}
\label{tab:ex4}
\begin{tabular}{lrr}
\toprule
Metric & Meal (60\,min) & IVGTT (2\,min) \\
\midrule
Peak plasma glucose ($\mgdl$) & $143.9$ & $263.4$ \\
Time-to-peak glucose (min) & $59.9$ & $2.0$ \\
Peak plasma insulin ($\uU$) & $30.5$ & $40.2$ \\
Time-to-peak insulin (min) & $13.8$ & $4.2$ \\
Glucose recovery (min) & $82.5$ & $53.9$ \\
\bottomrule
\end{tabular}
\end{table}

\textbf{Discussion.} The IVGTT dumps all 25\,g into the glucose pool within 2\,min, faster than any
clearance flux can respond, so glucose spikes to $263.4\mgdl$ almost immediately. The meal delivers the
identical mass at only $417$\,mg/min --- slow enough that the clearance fluxes ($F_3,F_4,F_5$) and the
suppression of hepatic production remove glucose as it arrives, giving a much lower, broader excursion
peaking at $143.9\mgdl$ near the end of the infusion. The insulin response mirrors the drive: the
gentler glucose rise opens the $W$ and $F_6$ secretion gates less strongly, so the meal insulin peak
($30.5\uU$) is blunted relative to the sharp IVGTT spike ($40.2\uU$), and then settles to a lower
sustained plateau as $r$ partially depletes and is refilled by synthesis $W$. This is why an oral meal
stresses glucose homeostasis far less than an equal-dose IV bolus, and why tolerance tests use a rapid
bolus to maximally probe the feedback system.

\begin{figure}[t]
\centering
\includegraphics[width=0.66\linewidth]{ex4_meal_vs_ivgtt.png}
\caption{Meal vs IVGTT: plasma glucose (overlay), plasma insulin (overlay), and the input rate
$I_g(t)$. Same 25\,g total, very different temporal spread.}
\label{fig:ex4}
\end{figure}

\subsection{Exercise 5: Hyperinsulinism and IV-glucose rescue}

\textbf{Question.} (a) Simulate hyperinsulinism (first guess: increase $a_6$) and terminate the patient
when plasma glucose falls below $20\mgdl$; report the fatal setting and time of death. (b) Resuscitate
with IV glucose: how much is needed?

\textbf{Approach (disclosed reformulation).} Simply raising $a_6$ does not kill the subject. At steady
state the pancreatic mass balance forces $F_6=W$, and the $W$ gate closes as glucose falls, so
endogenous insulin saturates ($\sim$40\uU) and glucose only dips to $\sim$78\mgdl\ before recovering.
This homeostasis is correct --- a normal pancreas cannot kill you --- so clinical insulin shock requires
an \emph{autonomous} source (insulinoma or overdose) that does not switch off when glucose drops. We
model this faithfully as an autonomous insulinoma with output
$I_{p,\mathrm{tumor}}=(M-1)\,F_{6,\mathrm{basal}}$ \si{\micro U}/min ($F_{6,\mathrm{basal}}=1.638\times10^{4}$
\si{\micro U}/min), where multiplier $M=1$ is the healthy subject. The native pancreas (the $W=F_6$
balance, all $\tanh$ shapes, all transfer rates) is unchanged; death is a terminal event on
$\bar{g}-20$.

\begin{table}[H]
\centering
\caption{Fatal-hyperinsulinism sweep. The first fatal multiplier is $4.5\times$; the demo scenario
uses $8\times$.}
\label{tab:ex5}
\begin{tabular}{rrrcr}
\toprule
$a_6$ mult & $I_{p,\mathrm{tumor}}$ (\si{\micro U}/min) & min $\bar{g}$ ($\mgdl$) & died & $t_{\mathrm{death}}$ (min) \\
\midrule
$1.0$ & $0$ & $91.50$ & no & --- \\
$3.0$ & $3.277\times10^{4}$ & $79.66$ & no & --- \\
$4.0$ & $4.915\times10^{4}$ & $36.65$ & no & --- \\
$4.5$ & $5.735\times10^{4}$ & $20.10$ & yes & $79.6$ \\
$6.0$ & $8.192\times10^{4}$ & $20.22$ & yes & $49.6$ \\
$8.0$ & $1.147\times10^{5}$ & $20.11$ & yes & $43.3$ \\
$10.0$ & $1.475\times10^{5}$ & $20.24$ & yes & $38.9$ \\
\bottomrule
\end{tabular}
\end{table}

\textbf{Results.} The lethal threshold is the $4.5\times$ multiplier; the demonstration scenario
($8\times$) kills the subject at $t=43.3$\,min (Fig.~\ref{fig:ex5}). For the rescue, a continuous IV
dextrose drip is given over a defined 240-min clinical window and the total dose found by bisection
(keeping $\bar{g}\ge20$ throughout): the \textbf{minimum rescue dose is $69.11$\,g} of glucose over
240\,min. Below this the nadir still crosses 20 and the patient dies; at or above it the glucose floor
holds at 20 and the patient survives the window.

\textbf{Limitation.} Because the modelled insulinoma is permanent, IV glucose is temporizing, not
curative (clinically the tumour must be removed) --- exactly the textbook's ``short-term treatment =
large IV glucose''. The minimum dose is therefore quoted for a defined window; a longer window needs
proportionally more glucose. Quoting an unconditional ``grams to survive forever'' would be undefined
for a permanent source.

\begin{figure}[H]
\centering
\begin{subfigure}{0.48\linewidth}
\includegraphics[width=\linewidth]{ex5_fatal_decline.png}
\caption{Fatal decline ($8\times$).}
\end{subfigure}\hfill
\begin{subfigure}{0.48\linewidth}
\includegraphics[width=\linewidth]{ex5_rescue.png}
\caption{Un-rescued death vs rescued survival.}
\end{subfigure}
\caption{Hyperinsulinism. Left: glucose driven through the 70 and 20\,mg/100\,ml lines as insulin rises
far above basal. Right: the $69.11$\,g/240\,min drip holds the glucose floor at 20.}
\label{fig:ex5}
\end{figure}

\subsection{Exercise 6: Insulin bolus, normal vs obese}

\textbf{Question.} Simulate rapid ingestion of insulin at $0.10$\,U/kg over 2\,min, observe the
momentary hypoglycemia, and ask whether the subject dies; repeat for an obese subject.

\textbf{Approach.} The dose is $0.10\,\text{U/kg}\times70\,\text{kg}=7\,\text{U}=7.0\times10^{6}$
\si{\micro U}, delivered as a rectangular $I_p(t)$ pulse ($3.5\times10^{6}$ \si{\micro U}/min) into the
plasma-insulin compartment over 2\,min. Each subject starts from its own pre-equilibrated basal; death
is a terminal event at $\bar{g}<20\mgdl$.

\begin{table}[H]
\centering
\caption{Insulin-bolus response. Both subjects reach the same insulin peak; only the glucose nadir
differs.}
\label{tab:ex6}
\begin{tabular}{lrrrl}
\toprule
Subject & glucose nadir ($\mgdl$) & $t_{\mathrm{nadir}}$ (min) & insulin peak ($\uU$) & verdict \\
\midrule
Normal & $20.08$ & $21.3$ & $1541$ & DIES ($t=21.4$\,min) \\
Obese & $59.83$ & $37.2$ & $1542$ & survives ($t_{\mathrm{rec}}=103.6$\,min) \\
\bottomrule
\end{tabular}
\end{table}

\textbf{Discussion.} Both subjects receive the identical 7\,U bolus and reach essentially the same
plasma-insulin peak ($\sim$1541\uU), yet their glucose nadirs differ sharply ($20.1$ vs $59.8\mgdl$).
The difference comes entirely from peripheral disposal ($F_4$): the obese subject freezes the
interstitial-insulin gate $H_4$ at $0.0012$ (Table~12.3) and shifts the glucose gate ($b_{42},c_{42}$),
so the same insulin flood removes far less glucose per minute. This is the model's representation of
insulin resistance, and here it is \emph{protective}: the normal subject's fully insulin-responsive
tissue drives glucose all the way through the $20\mgdl$ line and dies, while the resistant obese subject
bottoms out in the hypoglycemic-but-survivable range and recovers as the bolus clears and
counter-regulation restores glucose. The normal result sits right at the lethal boundary, but the
terminal event fires consistently across solvers and step sizes, so the verdict is robust within the
model.

\begin{figure}[t]
\centering
\includegraphics[width=0.66\linewidth]{ex6_insulin_bolus_normal_vs_obese.png}
\caption{Insulin bolus, normal vs obese. Identical insulin peaks, but the normal nadir touches
$20\mgdl$ (death) while the insulin-resistant obese subject bottoms out at $\sim$60\mgdl\ and recovers.}
\label{fig:ex6}
\end{figure}

\subsection{Exercise 7: A bowl of vanilla ice cream}

\textbf{Question.} Simulate the glucose, insulin and glucagon dynamics for a normal, diabetic and obese
subject each consuming an average bowl of vanilla ice cream.

\textbf{Approach and assumption.} An average bowl is taken as $\sim$1 cup ($\sim$130\,g) of regular
vanilla ice cream. USDA FoodData Central lists $\sim$23--24\,g carbohydrate per 100\,g (almost entirely
sugar), giving $\sim$31\,g; we adopt \textbf{30\,g of digestible carbohydrate, treated as 100\%
appearing as plasma glucose} (a conservative upper bound, since lactose and fat slow real absorption).
The bowl is eaten over 40\,min, modelled as a constant infusion $I_g=750$\,mg/min on $[0,40]$\,min --- a
first-order stand-in for the gastrointestinal $R_a(t)$. The normal subject starts from its clinical
basal, the obese subject is settled to its (normal-anchored) basal, and the diabetic uses the same
canonical configuration as Exercise~3 (its own elevated basal $140/7/75$).

\begin{table}[H]
\centering
\caption{Ice-cream meal (30\,g over 40\,min): per-subject peaks and recovery. Recovery is the first
return within 5\% of that subject's own basal.}
\label{tab:ex7}
\begin{tabular}{lrrrrrr}
\toprule
Subject & basal $g$ & peak $g$ & basal ins & peak ins & min glucagon & $g$ recovery (min) \\
 & ($\mgdl$) & ($\mgdl$) & ($\uU$) & ($\uU$) & (nU/ml) & \\
\midrule
Normal & $91.5$ & $185.7$ & $11.0$ & $36.9$ & $2.2$ & $77.3$ \\
Diabetic & $140.0$ & $252.8$ & $7.0$ & $9.4$ & $1.5$ & $126.4$ \\
Obese & $91.5$ & $206.8$ & $11.0$ & $35.5$ & $1.2$ & $113.8$ \\
\bottomrule
\end{tabular}
\end{table}

\textbf{Discussion.} A single identical meal produces three distinct responses. The \emph{normal}
subject rises to $\sim$186\mgdl\ then clears back toward $91.5$ with a clean first-phase insulin spike
($11\to37\uU$). The \emph{obese} subject reaches a slightly higher glucose peak ($\sim$207\mgdl) with a
broadly similar excursion shape from the same $91.5$ basal --- ``nearly normal'' glucose, as the
textbook states --- but reproduces the characteristic abnormal hump in the insulin-decay curve (a
secondary rise after the first-phase peak, driven by the steep $F_6$ gate and faster $m_{02}$
clearance). The \emph{diabetic} subject starts from an elevated fasting basal ($140\mgdl$), rises to the
highest peak ($\sim$253\mgdl) and recovers most slowly ($\sim$126\,min vs $\sim$77\,min for normal),
with honestly blunted insulin ($7.0\to9.4\uU$): impaired insulin-driven uptake ($F_2,H_4$ frozen low)
plus unsuppressed hepatic production (high $H_1$) gives a higher, slower excursion. In all three
subjects the meal-driven glucose and insulin rise suppresses glucagon (the $F_7$ gates are negatively
driven by both), so glucagon plunges from $\sim$75 toward $\sim$1--2\,nU/ml during the meal and then
recovers.

\begin{figure}[t]
\centering
\begin{subfigure}{0.32\linewidth}
\includegraphics[width=\linewidth]{ex7_icecream_glucose.png}
\caption{Plasma glucose.}
\end{subfigure}\hfill
\begin{subfigure}{0.32\linewidth}
\includegraphics[width=\linewidth]{ex7_icecream_insulin.png}
\caption{Plasma insulin.}
\end{subfigure}\hfill
\begin{subfigure}{0.32\linewidth}
\includegraphics[width=\linewidth]{ex7_icecream_glucagon.png}
\caption{Plasma glucagon.}
\end{subfigure}
\caption{Ice-cream meal, three subjects. Glucose: diabetic highest and slowest. Insulin: normal clean
spike, obese abnormal hump, diabetic blunted. Glucagon: suppressed in all three during the meal.}
\label{fig:ex7}
\end{figure}

\subsection{Part 1 summary}

The Cobelli model reproduces the qualitative IVGTT behaviour of Fig.~12.4 once an under-specified
parameter table is calibrated to a patient's basal operating point --- a step that turned out to be the
main practical obstacle and is documented in full rather than hidden. With that single calibrated core,
exercises 1--7 recover the expected physiology: CNS/RBC tissue takes most of an ingested glucose load,
the insulin hump arises from finite-pool secretion ($F_6\propto r$), a spread-out meal stresses
homeostasis far less than an equal-dose IV bolus, and diabetic, obese, hyperinsulinaemic and
insulin-overdose subjects each show their characteristic signature. Two reformulations were needed and
disclosed honestly: hyperinsulinism requires an autonomous (insulinoma) source because the intact
pancreas is self-limiting, and the diabetic subject must be calibrated to its own elevated operating
point because the normal basal is a strong attractor. None of these changes the model's feedback
structure.
